\section{All-order Stieltjes identities and Hurwitz primitives}
\label{dhr:sec:stieltjes}

\subsection{The complete jet formula}
Let $e_j=E^{(j)}(0)$, so $e_0=1$ and $e_1=\chi$. All these
coefficients are explicit finite polygamma polynomials. More precisely,
put
\begin{equation}\label{dhr:eq:ell}
 \ell_j=2\psi^{(j-1)}(1)-\psi^{(j-1)}(r)
        -\psi^{(j-1)}(a+r-1)-\psi^{(j-1)}(c+r-1).
\end{equation}
Then $\log E(u)=\sum_{j\geq1}\ell_j u^j/j!$ and
$e_j=\Bell_j(\ell_1,\ldots,\ell_j)$. In particular,
\[
 e_2=\chi^2+\ell_2,\qquad
 e_3=\chi^3+3\chi\ell_2+\ell_3.
\]
The derivatives $A_j$ are obtained in exactly the same way from the
finite Gamma product \eqref{dhr:eq:amplitude-E}.

\begin{theorem}[Every fixed Stieltjes order]\label{dhr:thm:jets}
For every integer $k\geq0$, the following single-bracket series is
ordinary and absolutely convergent:
\begin{align}
 &\sum_{n\geq0}\left[
 K_n^{(k)}(0)-\frac1{n+\alpha}
 \sum_{j=0}^k\binom{k}{j}A_j[-\log(n+\alpha)]^{k-j}
 \right]\notag\\
 &\qquad=\sum_{j=0}^k\binom{k}{j}A_j
 \left\{\frac{e_{k-j+1}}{k-j+1}
       -(-1)^{k-j}\gamma_{k-j}(\alpha)\right\}.
 \label{dhr:eq:all-jets}
\end{align}
The result is locally holomorphic for the parameters in
Theorem~\ref{dhr:thm:resonance}. The constants $A_j,e_j$ are finite
Gamma/polygamma expressions; no infinite coefficient extraction is
needed to compute a fixed row.
\end{theorem}
\begin{proof}
Differentiate the already normally convergent identity
\eqref{dhr:eq:D-master}, rather than the divergent unsubtracted
partition. The expansion at zero is
\[
 \frac{E(u)}u-\zeta(1+u,\alpha)
 =\sum_{h\geq0}\left\{\frac{e_{h+1}}{h+1}
               -(-1)^h\gamma_h(\alpha)\right\}\frac{u^h}{h!}.
\]
Leibniz's rule gives the right side of \eqref{dhr:eq:all-jets}.
Differentiating $(n+\alpha)^{-1-u}$ gives the logarithmic polynomial
on the left. The normal-convergence statement proved in
Theorem~\ref{dhr:thm:resonance} supplies absolute convergence of the
resulting bracket and justifies every interchange.
\end{proof}

The first three values $D_k(\alpha)=\partial_u^kD(0;\alpha)$ are
\begin{align}
 D_0(\alpha)&=A_0\{\chi-\gamma_0(\alpha)\},\notag\\
 D_1(\alpha)&=A_1\{\chi-\gamma_0(\alpha)\}
                   +A_0\{e_2/2+\gamma_1(\alpha)\},\notag\\
 D_2(\alpha)&=A_2\{\chi-\gamma_0(\alpha)\}
          +2A_1\{e_2/2+\gamma_1(\alpha)\}
          +A_0\{e_3/3-\gamma_2(\alpha)\}.
 \label{dhr:eq:first-jets}
\end{align}
The factor $1/(h+1)$ in these rows comes from dividing $E(u)-1$ by
$u$ before differentiation. Replacing it by one would give an
incorrect higher-order table.

\subsection{Derivatives in the subtraction point}
The parameter $\alpha$ specifies the subtraction grid. It does not
move the particles or change their original weights. Differentiating
this auxiliary family nevertheless yields useful exact identities
between the corresponding subtracted sums and the usual Hurwitz
order ladder.

\begin{proposition}[Hurwitz derivative and shift laws]\label{dhr:prop:alpha}
For every integer $h\geq1$,
\begin{equation}\label{dhr:eq:alpha-derivative}
 \partial_\alpha^hD(u;\alpha)
 =(-1)^{h+1}A(u)(1+u)_h\zeta(1+u+h,\alpha).
\end{equation}
Moreover,
\begin{equation}\label{dhr:eq:alpha-shift}
 D(u;\alpha+1)-D(u;\alpha)=A(u)\alpha^{-1-u}.
\end{equation}
Every fixed $u$-derivative of these equations is valid, including
at $u=0$. In particular
$\partial_\alpha^hD_0(\alpha)=A_0\psi^{(h)}(\alpha)$.
\end{proposition}
\begin{proof}
Use the classical identities
$\partial_\alpha^h\zeta(s,\alpha)=(-1)^h(s)_h\zeta(s+h,\alpha)$
and $\zeta(s,\alpha+1)=\zeta(s,\alpha)-\alpha^{-s}$
\cite{dlmf:hurwitz}. Their substitution in
\eqref{dhr:eq:D-master} proves both formulas. Joint holomorphy, already
established, permits the indicated mixed derivatives. The last formula
also follows immediately from $D_0=A_0(\chi+\psi(\alpha))$.
\end{proof}

\subsection{Anchored antiderivatives and log-Gamma}
An anchored primitive fixes the integration constant, which is especially
important when differentiating in the spectral variable. Let
$\Rea\alpha_0,\Rea\alpha>0$, with the path in that half-plane.

\begin{theorem}[Every anchored primitive of the spectral jets]
\label{dhr:thm:primitive}
The identity
\begin{equation}\label{dhr:eq:primitive-master}
 \int_{\alpha_0}^{\alpha}D(u;t)\diff t
 =\frac{A(u)}u\left\{(\alpha-\alpha_0)E(u)
                  +\zeta(u,\alpha)-\zeta(u,\alpha_0)\right\}
\end{equation}
has a removable singularity at $u=0$. For every $k\geq0$,
\begin{align}
 \int_{\alpha_0}^{\alpha}D_k(t)\diff t
 =\sum_{j=0}^k\binom{k}{j}\frac{A_j}{k-j+1}
 \big[& (\alpha-\alpha_0)e_{k-j+1}\notag\\
 &+\zeta^{(k-j+1)}(0,\alpha)
 -\zeta^{(k-j+1)}(0,\alpha_0)\big].
 \label{dhr:eq:primitive-jets}
\end{align}
In particular,
\begin{equation}\label{dhr:eq:gamma-primitive}
 \int_{\alpha_0}^{\alpha}D_0(t)\diff t
 =A_0\left\{(\alpha-\alpha_0)\chi
       +\log\Gamma(\alpha)-\log\Gamma(\alpha_0)\right\}.
\end{equation}
\end{theorem}
\begin{proof}
For $u\ne0$, integrate the identity
$\partial_t\zeta(u,t)=-u\zeta(1+u,t)$.
This gives \eqref{dhr:eq:primitive-master}. The numerator vanishes
at zero because $E(0)=1$ and $\zeta(0,t)=\tfrac12-t$.
Its divided Taylor series has $h$th derivative at zero
\[
 \frac{(\alpha-\alpha_0)e_{h+1}
       +\zeta^{(h+1)}(0,\alpha)-\zeta^{(h+1)}(0,\alpha_0)}{h+1}.
\]
Leibniz's rule proves \eqref{dhr:eq:primitive-jets}. Finally,
$\zeta'(0,t)=\log\Gamma(t)-\tfrac12\log(2\pi)$
\cite{dlmf:hurwitz} gives \eqref{dhr:eq:gamma-primitive}.
\end{proof}

This theorem transports the classical Stieltjes primitive to a
many-variable resonant harmonic sum. It does not claim a new formula
for $\int\gamma_k(t)\diff t$ in isolation, nor does it integrate an
arbitrary many-variable weight with respect to a particle coordinate.
The distinction identifies precisely which antiderivative has been
established.
